\section{Experimental Design and Expected Results}
\subsection{Setup}
Scenarios use fleet sizes $n\in\{5,10,20,50\}$, $m$ non-cooperative intruders, static
obstacles of varying density, and communication latencies
$\delta\in\{0,0.1,0.5,1.0\}$~s. The GNN encoder is trained jointly across
$n \in \{3, 5, 8\}$ using the parallel training protocol of Section~\ref{subsec:gnn}
and evaluated zero-shot at $n \in \{10, 20, 50\}$. Loss-of-separation events are
deliberately induced during training to expose the conflict-resolution policy to the
collision regime; a separate, more realistic scenario is used to evaluate the handoff
between the PID controller and the MARL policy. Baselines are QMIX~\cite{qmix},
MAPPO~\cite{mappo}, GCBF~\cite{gcbf}, and a local-only ablation in which the
conflict-resolution policy is removed and the PID controller operates alone. Metrics are
mission success rate, loss-of-separation rate, path efficiency, intervention frequency,
and emergence measures: heading-direction entropy, trajectory mutual information, and
deadlock count.

% TODO(authors): the implementation now uses BlueSky / BlueSky-Gym (PettingZoo branch)
% \cite{bluesky,bluesky-gym}. Decide whether BlueSky replaces the "2D navigation
% benchmark" below or is added as a third environment, and update this paragraph.
Validation is conducted in two environments. The first is a standard 2D multi-agent
navigation benchmark used widely in social force and robot motion planning research,
providing a controlled setting in which fleet size can be varied systematically. The
second is a custom 3D UAV simulation environment developed by our research team, which
models UAV dynamics with six-degree-of-freedom kinematics, realistic airspace geometry, and stochastic agent entry and exit
events that mimic AAM operational patterns.

\subsection{Hypothesized Outcomes}
The central hypothesis is that the proposed architecture retains a high mission success
rate as fleet size grows, because the GNN encoder produces a fixed-size latent vector
independent of $n$ and the downstream Q-network is consequently decoupled from fleet size,
whereas CTDE baselines with fixed-dimensional mixing or critic networks cannot be evaluated
at a fleet size other than the one they were trained on without architectural change.
We further expect that removing the conflict-resolution policy (the local-only ablation)
degrades success in high-density scenarios, since the PID controller alone has no mechanism
for fleet-level coordination. Table~\ref{tab:expected} states these expectations as a
qualitative prediction; the entries are hypotheses to be tested, not measured results.


\begin{table}[hbt!]
\caption{\label{tab:expected} Hypothesized zero-shot mission success after training at
$n=5$. Entries are predicted outcomes to be tested in the final manuscript, not measured
results.}
\centering
\begin{tabular}{lcccc}
\hline
Method & $n=5$ & $n=10$ & $n=20$ & $n=50$ \\
\hline
Proposed (hybrid, GNN encoder) & High & High & High & Moderate--High \\
QMIX / MAPPO & High & \multicolumn{3}{c}{Not transferable without retraining} \\
Local-only ablation (PID only) & Moderate & Moderate & Low & Low \\
\hline
\end{tabular}
\end{table}


\subsection{Planned Evaluation and Deliverables}
The final paper will report: (i) a quantitative scaling curve of mission success rate
against fleet size $n$, with the empirically fitted degradation exponent for the proposed
method and for the CTDE baselines; (ii) the effect of communication latency $\delta$ on
success rate, compared against the theoretical degradation characterization; (iii)
ablations isolating the conflict-resolution policy and the PID input, including a
comparison of sample efficiency against end-to-end MARL; and (iv) emergence analysis
quantifying lane formation through heading-direction entropy, leader--follower structure
through trajectory mutual information, and deadlock frequency relative to baselines.
The theoretical claims stated in Section~\ref{sec:intro} will be accompanied by formal
statements and proofs. The simulation environment, the PID local controller, the
decentralized broadcast channel, and the GNN encoder with its downstream Q-network are
under active implementation, which supports completion of this evaluation by the manuscript
deadline.

